% cap1.tex -- capitulo unico do projeto de teste (fixture) da Fase 0. Reune,
% de forma compacta, um caso conhecido de cada checagem objetiva do pipeline.
\section{Capitulo Um}

O presente trabalho utiliza uma metodologia experimental para avaliar o desempenho do sistema proposto.

O sistema utiliza uma API para comunicacao entre os modulos internos, permitindo a troca de mensagens em tempo real.

O projeto foi implementado com o auxilio de um framework de desenvolvimento agil, o que acelerou o ciclo de testes.

Trata-se de uma abordagem inovadora para o processamento digital de sinais, conforme descrito na literatura correlata \cite{smith2020}.

Este trabalho descreve, de maneira detalhada e abrangente, o desenvolvimento de um sistema completo de aquisicao, processamento, armazenamento e transmissao de dados provenientes de sensores distribuidos ao longo de uma rede de comunicacao sem fio, considerando aspectos de desempenho, confiabilidade, seguranca da informacao e eficiencia energetica durante toda a operacao continua do sistema proposto pelos autores deste trabalho.

Conforme mostra a Figura~\ref{fig:naoexiste}, o sistema apresenta melhorias significativas em relacao a versao anterior, um resultado tambem relatado por outro autor \cite{chaveFantasma}.

\begin{figure}
  \centering
  \includegraphics{figuras/ausente.png}
  \caption{Grafico de desempenho do sistema}
  \label{fig:grafico}
\end{figure}

\begin{figure}
  \centering
  Uma segunda figura, apresentada apenas como texto ilustrativo, sem legenda e sem rotulo associado.
\end{figure}

Mais adiante, convem esclarecer que a Interface de Programacao de Aplicacoes (API) e o mecanismo central descrito acima.

O experimeto de validacao foi conduzido em ambiente controlado de laboratorio.

O presente trabalho utiliza uma metodologia experimental para avaliar o desempenho do sistema proposto.
